\documentclass[Total.tex]{subfiles}

\begin{document}
	\chapter{Differentiable Manifolds}
	
		\section{Differerentiable Manifolds}
			\subsection{Pseudo Group Axioms}
				
				\begin{Def}
					A pseudogroup of transformations on a topological space $S$ is a set $\Gamma$ of transformations satisfying the following axioms:
					
					\begin{itemize}
						\item Each $f\in\Gamma$ is a homeomorphism of an open set of $S$ onto another open set of $S$.
						
						\begin{gather*}
							f_{\cong}:U\rightarrow V;U,V\in Open(S)
						\end{gather*}
						\item If $f\in \Gamma$, then the restriction of $f$ to an arbitrary open subset of domain of $f$ is in $\Gamma$.(And it is still an homeomorphism)
						
						\begin{gather*}
							f\in \Gamma\Rightarrow f|U\in \Gamma
						\end{gather*}
						
						\item Let $U=\cup U_i$. Then, 
						
						\begin{gather*}
							f|U_i\in \Gamma\Rightarrow f\in \Gamma
						\end{gather*} 
						\item For every open $U$,
						
						\begin{gather*}
							id_U\in \Gamma
						\end{gather*}
						\item 
						
						\begin{gather*}
							f\in \Gamma\Rightarrow f^{-1}\in \Gamma
						\end{gather*}
						\item $f\in \Gamma$ is a homeomorphism of $U$ onto $V$, $f'\in \Gamma$ is a homeomorphism of $U'$ onto $V'$, then 
						
						\begin{gather*}
							\Gamma\ni f'\circ f:f^{-1}(V\cap U')\rightarrow f'(V\cap U')
						\end{gather*}
					\end{itemize}
				\end{Def}
				
				\begin{example}
					(Pseudogroup $\Gamma^r(\RR^n)$) Define:
					
					\begin{gather*}
						\Gamma^r(\RR^n):=\{f\in C^r(U,V)|U,V\in Open(\RR^N)\}
					\end{gather*}
					
					There is an inclusion
					
					\begin{gather*}
						\Gamma^{r+1}(\RR^n)\hookrightarrow \Gamma^{r}(\RR^n)
					\end{gather*}
				\end{example}
				
				\begin{example}
					Furthermore, Psudogroup
					
					\begin{gather*}
						\mathcal{O}(\Omega)
					\end{gather*}
				\end{example}
				
				\subsection{Atlas}
				
				\begin{Def}
					An \textbf{atlas} of a topological space $M$ is a family of \textbf{charts} $(U_i,\varphi_i)$ which satisfies:
					
					\begin{itemize}
						\item Each $U_i$ is an open set of $M$, and 
						
						\begin{gather*}
							\cup U_i=M
						\end{gather*}
						\item $\varphi_i:U_i\rightarrow S$ is a homeomorphism onto its image;
						\item If $U_i\cap U_j\not=\varnothing$, then
						
						\begin{gather*}
							\varphi_j\circ \varphi_i^{-1}:\varphi_i(U_i\cap U_j)\rightarrow \varphi_j(U_i\cap U_j)
						\end{gather*}
						
						is an element of $\Gamma$.
					\end{itemize}
				\end{Def}
				
				\begin{Def}
					An \textbf{compelete atlas} of $M$ compatible with $\Gamma$ is an atlas of $M$ compatible with $\Gamma$.
				\end{Def}
				
				Then, every atlas is contained in a unique compelete atlas. 
				
				\textbf{Remark} Suppose $A,A'$ be two atlases. Then, if they are \textbf{compatible}
				
				\begin{gather*}
					内容...
				\end{gather*}
				
				\subsection{Differentiable Manifold}
				
				\begin{Def}
					内容...
				\end{Def}
				
				\subsection{Differentiable Mapping}
				
				\subsection{Differentiable Curve, Tangewnt Vectors}
					
					\subsubsection{Differentiable Curve}
					
						\begin{Def}
							内容...
						\end{Def}
						
					\subsubsection{Tangent Vectors}
						
						\begin{Def}
							Let $\mathfrak{F}(p),p\in M$ be the equivalence class of differentiable functions at $p$,
							
							\begin{gather*}
								\mathfrak{F}(p):=\{(U,f)|f\in C^1(U,\RR),p\in U\subseteq X\mbox{ open }\}/\sim\qquad (U,f)\sim (V,g)\Leftrightarrow f|U\cap V=g|U\cap V
							\end{gather*}
							
							$\mathfrak{p}$ is an $\RR$-linear space.
						\end{Def}
						
						Then, let $x:[a,b]\rightarrow M$ be a $C^1$-differentiable function where
						
						\begin{gather*}
							x(t_0)=p
						\end{gather*}
						
						Then, define 
						
						\begin{gather*}
							X:\mathfrak{F}(p)\rightarrow \RR\\
							[(U,f)]\mapsto (\frac{df(x(t))}{dt})|_{t_0} 
						\end{gather*}
						
						In other word, $Xf$ is derivative of $f$ in the direction of the curve $x(t)$ at $t=t_0$. 
						
						The function $X$ satisfies the following properties:
						
						\begin{itemize}
							\item 
							\item 
						\end{itemize}
						
					\subsubsection{Lie Bracket}
					
		\section{Tensor Algebra, Tensor Fields}
			Given in Smooth manifold.
		\section{Lie Groups}
			Given in Lie theory.
		\section{Fibre Bundles}
			\subsection{Differentiable Principle Fibre Bundle}
			
				\begin{Def}
					A \textbf{differentiable principle fibre bundle} over $M$ with group $G$ consists of $P(M,G,\pi)$, where:
					
					\begin{itemize}
						\item $G$ acts freely on $P$ on the right:
						\begin{gather*}
							P\times G\rightarrow P\qquad(u,a)\mapsto ua=R_au\in P
						\end{gather*}
						\item $M$ is quotient space of $P$, and the canonical projection
						\begin{gather*}
							\pi:P\rightarrow M
						\end{gather*}
						is differentiable;
						
						\item $P$ is locally trivial:
					\end{itemize}
				\end{Def}
				
				We call:
				
				\begin{itemize}
					\item 
				\end{itemize}
				
				\subsubsection{Fibre of Differentiable Principle Fibre Bundle}
				
				\subsection{Fundamental Vector Field} 
				\begin{Def}
					Given a principle fiber bundle $P(M,G,\pi)$, the action of $G$ on $P$ induces a homomorphism $\sigma$ of the Lie algebra $\mathfrak{g}$ of $G$ into the Lie algebra $\mathfrak{X}(P)$ of vector fields on $P$.
					\begin{gather*}
						\mathfrak{g}=L(G)\rightarrow \mathfrak{X}(P) 
					\end{gather*}
				\end{Def}
				Constructions:
				\begin{itemize}
					\item Exponential Maps:
					\item Consider the action 
					\begin{gather*}
						\Phi_u:G\rightarrow P\quad g\mapsto u\cdot g
					\end{gather*}
					Then,
					\begin{gather*}
						\sigma:\mathfrak{g}=T_e G\rightarrow \mathfrak{X}(P	) \sigma(A)_u=(d\Phi_u)_e(A)
					\end{gather*}
				\end{itemize}
				Take the following notation
				\begin{gather*}
					A^*:=\sigma A
				\end{gather*}
				Furthermore, we have
				\begin{proposition}
					\begin{gather*}
						\mathfrak{g}\rightarrow T_u P\qquad A\mapsto (A^*)_u
					\end{gather*}
					is linear isomorphism.
				\end{proposition}
				\begin{proof}
					\begin{itemize}
						\item 
						\item 
						\item 
					\end{itemize}
				\end{proof}
				
				\begin{proposition}
					Let $A^*$ be fundmanetal vector field, then there exists correspondence
					\begin{gather*}
						A\in \mathfrak{g}\Rightarrow A^*\in \mathfrak{X}(P)\\
						(ad(a^{-1}))A\in \mathfrak{g}\Rightarrow (R_a)_*A^*
					\end{gather*}
				\end{proposition}
				\begin{proof}
					内容...
				\end{proof}
				
				The mapping $\sigma$ is injective, but not surjective. But we could define a Lie-subalgebra of $\mathfrak{X}(P)$, so that $\sigma:L(G)\rightarrow \mathfrak{X}(P)$ is bijective.
				
				\begin{Def}
					内容...
				\end{Def}
				
				\begin{proposition}
					Let $M$ be a manifold. $\{U_\alpha\}$ is an open covering of $M$ and $G$ a Lie group. 
					
					Given mapping 
				\end{proposition}
				\begin{proof}
					内容...
				\end{proof}
	
	\chapter{Theory of Connections}
		\section{Connections in a Principle Vector Bundle}
			\textbf{Remark} Firstsly, the fibre through $u$ is 
			\begin{gather*}
				E_{\pi(u)}
			\end{gather*}
			
			\begin{Def}
				Let $(P,M,G,\pi)$ be a principle fibre bundle over a manifold $M$. For each $u\in P$, $T_u(P)$ be the tangent space of $P$ at $u$, and $G_u$ be the subspace of $T_u(P)$ consisting of vectors tangent to fibre through $u$. 
				\begin{gather*}
					G_u:=T_u(E_{\pi(u)})=Ker(\pi_u)
				\end{gather*}
				\begin{itemize}
					\item $T_u(P)=G_u\oplus Q_u$; $G_u$ is the verticle subspace, $Q_u$ is the horizontal subspace;
					
					Take the following notation $vX,hX$. 
					\item $Q_{ua}=(R_a)^*Q_u;u\in P,a\in G$
					\item $Q$ depends differentiably on $u$;
				\end{itemize}
				A \textbf{connection} is a distribution $u\mapsto Q_u$. 
			\end{Def}
			
			\textbf{Remark 1} The tangent space is the formel tangent space, not rely on the open subset and the tangent line via limit. 
			
			\subsection{Connection Form}
			
			\begin{lemma}
				$\sigma_u:\mathfrak{g}\rightarrow T_u(P)$ is injective, $u\in P$,with induced linear isomorphism
				\begin{gather*}
					\sigma_u:\mathfrak{g}\rightarrow V_u(P)\qquad A\mapsto (A^*)_u
				\end{gather*}
			\end{lemma}
			\begin{proof}
				内容...
			\end{proof}
			
			\begin{Def}
				Given a connection $\Gamma$ in $P$. We define the $1$-form $P$ with the value in Lie algebra of $G$ as follows: $\omega\in \Omega^1(P,\mathfrak{g})$,
				\begin{gather*}
					\omega_p:T_u(P)\rightarrow \mathfrak{g}\qquad X\mapsto A
				\end{gather*}
				
				such that $(A^*)_u=vX$.
			\end{Def}
			
			\textbf{Remark 1} $\mathfrak{g}\rightarrow G_u,A\mapsto (A^*)_u$ is a linear isomorphism.
			
			\begin{proposition}
				The connection form $\omega$ of a connection satisfies the following conditions:
				\begin{itemize}
					\item $\omega(A^*)=A$;
					\item $(R_a)^*=ad(a^{-1})\omega$, that is
					\begin{gather*}
						内容...
					\end{gather*}
					for every $a\in G$ and every vector field $X$ on $P$. 
				\end{itemize}
			\end{proposition}
			\begin{proof}
				内容...
			\end{proof}
			
			\begin{proposition}
				内容...
			\end{proposition}
			\begin{proof}
				内容...
			\end{proof}
			
			\begin{proposition}
				内容...
			\end{proposition}
		
		\section{Existence and Extension of Connections}
			
			\begin{theorem}
				内容...
			\end{theorem}
			\begin{proof}
				内容...
			\end{proof}
		\section{Parallelism}
			\subsection{Horizontal Lift}
				
				\begin{Def}
					Let $\tau=x(t)$ be a piecewise differentiable curve of $C^1$ in $M$. A \textbf{horizontal lift} or \textbf{lift} is a \textbf{horizontal curve}
					\begin{gather*}
						\tau^*=u(t)
					\end{gather*}
					such that $x(t)=\pi(u(t))$. 
					
					Horizontal: Piecewise differentiable, the tangent vectors are horizontal.
				\end{Def}
				
				\begin{proposition}
					(Existence, Uniqueness)
				\end{proposition}
				\begin{proof}
					内容...
				\end{proof}
				
			\subsection{Parallel Displacement}
		
		\section{Holonomy Groups}
		
			Another Reference: Riemannian Holonomy Groups
			and Calibrated Geometry by Dominic D. Joyce
			\begin{Def}
				For point $x\in M$, denote
			\end{Def}
		\section{Curvature Form, and Structure Equation}
				\begin{Def}
					Let $P(M,G,\pi)$ be a PFB and representation on $G$ on $V$, with $\rho(a)\in End(V),a\in G$. $\rho(ab)=\rho(a)\rho(b)$. 
					
					A \textbf{pseudotensorial form} of degree on $P$ of type $(\rho,V)$ is a $V$-valued $r$-form $\varphi$ such that
					\begin{gather*}
						R_a^*\varphi=\rho(a^{-1})\cdot \varphi;a\in G
					\end{gather*} 
					
					Such a form is called \textbf{tensorial form} if it is horizontal in the sense of
					\begin{gather*}
						\varphi(X_1,\dots,X_i,\dots,X_r)=0;X_i\in V_i P
					\end{gather*}
				\end{Def}
				
				\begin{proposition}
					\begin{itemize}
						\item 
						\item 
						\item 
					\end{itemize}
				\end{proposition}
				
				\begin{proof}
					内容...
				\end{proof}
				
				Especially:
				\begin{Def}
					\begin{itemize}
						\item The connection form $\omega$ is a pseudotensorial 1-form of type $ad(G)$;
						\item $D\omega$  is called \textbf{tensorial 2-form} of type $ad(G)$, denoted by $\Omega$, is called the \textbf{curvature form} of $\omega$. 
					\end{itemize}
				\end{Def}
				
				\subsection{Structure Equation}
				
				\begin{theorem}
					(Structure Equation,E.Cartan) Let $\omega$ be a connection form and $\Omega$ be its curvature form, then
					\begin{gather*}
						\Omega(X,Y)=d\omega(X,Y)+\frac{1}{2}[\omega(X),\omega(Y)]
					\end{gather*}
				\end{theorem}
				\begin{proof}
					内容...
				\end{proof}
				
				\begin{proposition}
					内容...
				\end{proposition}
				\begin{proof}
					内容...
				\end{proof}
				
				\subsection{Bianchi's Identity}
				
				\begin{theorem}
					(Bianchi)
					\begin{gather*}
						D\Omega=0
					\end{gather*}
				\end{theorem}
				\begin{proof}
					内容...
				\end{proof}
		\section{Mappings of Connections}
		\section{Reduction Theorem}
		\section{Holonomy Theorem}
		\section{Flat Connections}
		\section{Local, and Infinitesimal Holonomy Groups}
		\section{Invariant Connections}
	\chapter{Linear Affine Connections}
		\section{Connections in Vector Bundles}
			\subsection{Fibre Metrics}
		\section{Linear Connections}
	
	\chapter{Curvature and Space Forms}
	
	\chapter{Transformations}
	
	\chapter{Submanifolds}
	
	\chapter{Variation of Length Integral}
	
	\chapter{Homogeneous Spaces}
	
	\chapter{Symmetric Spaces}
\end{document}